\documentclass{beamer}
\usepackage{hyperref}
\usepackage{fontspec}
\usepackage{braket}
\usepackage{natbib}
\usepackage{graphicx}
\usepackage{tikz}
\usetikzlibrary{arrows.meta, positioning, decorations.pathreplacing, calc}

\usepackage{lipsum}



% other packages
\usepackage{latexsym,amsmath,xcolor,multicol,booktabs,calligra}
\usepackage{graphicx,pstricks,listings,stackengine}

% dummy text; remove it when working on this template
% \usepackage{lipsum}

% this line disabled navigation toolbar
\setbeamertemplate{navigation symbols}{}

\usepackage{ariana}
\newcommand{\pinkaccent}[1]{\textcolor{barbiepink}{#1}}


\setbeamercolor{alerted text}{fg=barbiepink}
\newcommand{\muted}[1]{{\color{dimgray}#1}}

\title    {The Revelation Principle in Lean 4}
\subtitle {A Simulation theorem for mediated protocols}
\author   {Ariana Tang}
\institute{MATH 36610}
\date     {November 2026}

\begin{document}

% ======================================================
\begin{frame}
  \titlepage
\end{frame}


% ======================================================
\begin{frame}{The setup}
  \vspace{0.6em}
  \centering
  \begin{tikzpicture}[
      node distance=0.6cm,
      box/.style={draw=deeptextgray,rounded corners=2pt,inner sep=5pt,font=\small},
      lbl/.style={font=\small},
      sym/.style={font=\small\color{dimgray}},
      ar/.style={-{Stealth[length=5pt]},draw=deeptextgray}]
    \node[lbl] (t) {type};
    \node[box,right=of t] (p1) {player};
    \node[box,right=1.55cm of p1] (med) {mediator};
    \node[box,right=1.55cm of med] (p2) {player};
    \node[lbl,right=of p2] (a) {action};
    \draw[ar] (t) -- (p1);
    \draw[ar] (p1) -- node[above,font=\small]{message} (med);
    \draw[ar] (med) -- node[above,font=\small]{signal} (p2);
    \draw[ar] (p2) -- (a);
    \node[sym,below=0.35cm of t] {$T$};
    \node[sym,below=0.35cm of a] {$A$};
    \node[sym] at ($(p1)!0.5!(med)+(0,-0.6)$) {$M$};
    \node[sym] at ($(med)!0.5!(p2)+(0,-0.6)$) {$R$};
  \end{tikzpicture}

  \vspace{1.4em}
  \raggedright
  \begin{itemize}\setlength\itemsep{0.5em}
    \item Each player privately learns a type, talks to the mediator, privately hears back, then acts.
    \item \textbf{Payoffs:} types $t$ are drawn from a common prior $p$; player $i$ gets $u_i(a, t)$.
    \item \textbf{Strategy:} what to say given my type; what to do given what I've seen.
    \item \textbf{Equilibrium:} no player raises their expected payoff by changing their own strategy alone.
  \end{itemize}
\end{frame}


% ======================================================
\begin{frame}[fragile]{The theorem}
  Any equilibrium outcome, with \emph{any} message and signal sets, is the outcome of
  \alert{truthful, obedient} play with a \alert{direct} mediator: $M = T$, $R = A$.

  \vspace{0.6em}
\begin{minted}{lean4}
theorem revelation_principle_iff (G : BayesGame T A)
    (μ : (∀ j, T j) → FinDist (∀ j, A j)) :
    IsCommunicationEqOutcome G μ
      ↔ ∃ (d' : DirectDevice T A),
          IsCommunicationEq G d' truthfulProfile
          ∧ ∀ t, outcome d' truthfulProfile t = μ t
\end{minted}

  \vspace{0.2em}
  \muted{where the left side quantifies over all protocols:}
\begin{minted}{lean4}
∃ (M R : ι → Type) (_ : ∀ i, Fintype (M i)) (_ : ∀ i, Fintype (R i))
  (d : Device M R) (π : MixStrategyProfile T A M R), ...
\end{minted}
\end{frame}

% ======================================================
\begin{frame}[fragile]{Key construction}
  \vspace{0.8em}
\begin{minted}{lean4}
def canonical (d : Device M R) (π : MixStrategyProfile T A M R)
    : DirectDevice T A :=
  Device.mk (outcome d π)
\end{minted}

  \vspace{1.2em}
  \begin{itemize}\setlength\itemsep{0.7em}
    \item The new mediator \emph{is} the old outcome function: it runs the old protocol internally, players included.
    \item Truthful play reproduces the outcome: \texttt{pi\_pure}, \texttt{pure\_bind}, \texttt{map\_id'}.
    \item Any deviation against the new mediator can be \alert{simulated} by a deviation against the old one, so it can't pay.
  \end{itemize}
\end{frame}

% ======================================================
\begin{frame}[fragile]{Design: probability as a monad}
\begin{minted}{lean4}
structure FinDist (α : Type) [Fintype α] where
  w       : α → ℝ
  nonneg  : ∀ a, 0 ≤ w a
  sum_one : ∑ a, w a = 1
\end{minted}

\begin{minted}{lean4}
def outcome (d : Device M R) (π : MixStrategyProfile T A M R)
    (t : ∀ i, T i) :=
  FinDist.bind (FinDist.pi π) (fun s => play d s t)
\end{minted}

  \vspace{0.4em}
  \begin{itemize}\setlength\itemsep{0.6em}
    \item Hand-rolled finite distributions with \texttt{pure}, \texttt{bind}, \texttt{map}, \texttt{pi}; the game is Kleisli composition.
    \item Player-dependent sets are dependent types: \mbox{\texttt{T : ι → Type}}, profiles \mbox{\texttt{∀ i, T i}}.
  \end{itemize}
\end{frame}

% ======================================================
\begin{frame}[fragile]{What the paper hid}
  On paper, ``$i$ deviates, others keep playing'' is one symbol: $(\sigma_i', \sigma_{-i})$.

  \vspace{0.4em}
  In Lean, it is \texttt{Function.update π i ρ}, and pushing it through the product needs
\begin{minted}{lean4}
lemma pi_update_bind (μ : ∀ i, FinDist (X i)) (i : ι)
    (ν : FinDist (X i)) (F : (∀ j, X j) → FinDist γ) :
    (§$\bigotimes$§ (Function.update μ i ν)) >>= F
      = ν >>= (fun x => (§$\bigotimes$§ μ) >>= (fun s => F (Function.update s i x)))
\end{minted}
  \muted{plus four bookkeeping lemmas splitting $s$ into $(s_i, s_{-i})$ via \texttt{Equiv.piSplitAt}.}

  \vspace{0.8em}
  Meanwhile the equilibrium argument itself: \alert{about 5 lines}.\\
  \muted{Easy on paper and hard in Lean are inverted.}
\end{frame}

% ======================================================
\begin{frame}{Status and next steps}
  \vspace{0.6em}
  I don't know yet.
  % \begin{columns}[T,onlytextwidth]
  %   \column{0.48\textwidth}
  %   \textbf{Done} \muted{\small (no \texttt{sorry})}
  %   \begin{itemize}\setlength\itemsep{0.5em}
  %     \item \texttt{FinDist} and its monad laws
  %     \item Bayesian games, devices, equilibrium
  %     \item Truthful play reproduces the outcome
  %     \item Deviations are simulated in the original protocol
  %     \item Revelation principle, both directions
  %   \end{itemize}

  %   \column{0.48\textwidth}
  %   \textbf{Next}
  %   \begin{itemize}\setlength\itemsep{0.5em}
  %     \item Outcome set as a convex polytope cut out by linear incentive constraints
  %     \item Computable version ($\mathbb{Q}$ weights) to \texttt{\#eval} examples
  %     \item Contribute to EconLib
  %   \end{itemize}
  % \end{columns}
\end{frame}

\end{document}